\begin{abstract}
Quantum machine-learning components are often reported to improve classical
pipelines by less than the pipelines' own run-to-run variation. We study
whether quantum components help exemplar-free class-incremental learning with
a pre-trained model, under a protocol that resolves a few tenths of a point:
paired runs, all inference heads evaluated on one network, decision rules
fixed in advance, and a classical twin with the same inputs for every quantum
component. We test seven
placements: a projected quantum kernel as a steering metric, a separation
regulariser and a feature map; the Holevo quantity and a measurement-based
loss as training objectives; and a pretty-good measurement over class density
matrices, alone and fused with the classifier. Only one twin comparison has an
interval above zero, at $+0.04$ points. A read-only probe shows that the separation
kernel loses class information in its 8-qubit projection; repairing this lowers
accuracy. We prove that circuits after the encoding and state shrinkage
cannot change the measurement, and show across 61 stored runs that every
training-time change trades new-task for old-task accuracy. A calibrated fusion
of the classifier with a ridge score on the stored class statistics raises final
accuracy on four benchmarks by $0.38$ to $0.67$ points; the quantum score does
the same.
\end{abstract}
